\chapter{Implicazioni organizzative della Direttiva NIS2}
\label{cap:nis2}

% 1.1 %%%%%%%%%%%%%%%%%%%%%%%%%%%%%%%%%%%%%%%%%%%%%%%%%%%%%%%%%%%%%%%%%%%%%%%%
\section{Confronto tra le Direttive NIS e NIS2}
La \textit{Direttiva (UE) 2016/1148}\footcite{nis1-directive}, nota come direttiva NIS, ha rappresentato il primo tentativo del
legislatore europeo di armonizzare la cybersicurezza all'interno dell'Unione.

La direttiva mirava a innalzare il livello di sicurezza delle reti e dei sistemi informativi, istituendo 
un quadro comune volto a garantire maggiore uniformità tra gli Stati membri sugli standard minimi di 
cybersicurezza e di resilienza alle minacce. La sua adozione ha avuto implicazioni rilevanti sul 
piano organizzativo, in quanto ha riconosciuto formalmente la rilevanza strategica delle infrastrutture digitali 
per il mercato interno e ha introdotto, per la prima volta a livello europeo, un approccio 
uniforme di contrasto alle minacce informatiche basato sulla gestione del rischio.

Nonostante il suo ruolo pionieristico, la prima direttiva NIS ha mostrato evidenti limiti
che ne hanno ridotto l'efficacia e l'impatto rispetto alle aspettative. In primo luogo, il recepimento a livello
dei singoli Stati membri si è dimostrato fortemente disomogeneo, producendo un quadro frammentato
a livello europeo\footcite[p. 5]{vandezande}. Inoltre, la normativa presentava anche un problema intrinseco: 
il suo campo di applicazione era troppo ristretto, includendo soltanto alcuni settori considerati critici 
(energia, trasporti, bancario, finanziario, sanità, acqua potabile e infrastrutture digitali). Il limitato
raggio di azione ha rappresentato un grosso limite nel tentativo di diffondere un approccio uniforme, escludendo
importanti aree anch'esse sensibili alle minacce informatiche\footcite[p. 209]{teichmann}.
A ciò si aggiungeva la debolezza dei meccanismi di supervisione e sanzione, lasciati in larga misura alla 
discrezionalità nazionale, con conseguenti disparità nel livello di implementazione e nella 
deterrenza effettiva\footcite[p. 6]{vandezande}.

In risposta a tali criticità, l’Unione europea ha adottato la \textit{Direttiva (UE) 2022/2555} 
(NIS2), entrata in vigore a gennaio 2023 e destinata ad abrogare la NIS dal 18 ottobre 
2024. NIS2 rafforza la portata delle misure previste dal testo precedente ampliando i settori di interesse,
introducendo la distinzione tra soggetti essenziali e importanti, imponendo requisiti più 
specifici di gestione del rischio e notificazione degli incidenti, nonché un regime sanzionatorio 
più severo e uniforme a livello europeo. Anche l'obiettivo dichiarato evolve verso l'attenzione alla cybersicurezza 
in senso più ampio: non si limita alla protezione della rete e dei sistemi informatici, ma anche degli utenti di
tali sistemi e di tutti coloro interessati dalle minacce\footcite[p. 5]{vandezande}. In questo
senso, NIS2 si pone come il primo passo verso un framework europeo per la cybersicurezza, 
dimostrando una forte vocazione operativa nelle misure prescritte.



% 1.2 %%%%%%%%%%%%%%%%%%%%%%%%%%%%%%%%%%%%%%%%%%%%%%%%%%%%%%%%%%%%%%%%%%%%%%%%
\section{Misure principali della Direttiva NIS2}

\subsection{Ambiti e soggetti coinvolti}
In linea generale, rientrano nel campo di applicazione della direttiva tutte le imprese di 
medie e grandi dimensioni che operano nei settori individuati dal legislatore come rilevanti per la resilienza 
dell’Unione\footcite[art. 2]{nis2-directive}. Le micro e piccole imprese non sono dunque contemplate, salvo i casi in cui forniscano servizi di 
natura critica, svolgano un ruolo unico o rappresentino un monopolio, forniscano servizi ad un soggetto interessato
o siano specificatamente designate dal rispettivo Stato membro.

All’interno di tutti gli attori così individuati, la normativa introduce una distinzione in due categorie:
\begin{itemize}
    \item \textbf{Soggetti essenziali}, ossia gli operatori di settori la cui interruzione avrebbe conseguenze sistemiche 
        per la società e l’economia (tra cui energia, trasporti, sanità, infrastrutture digitali)\footcite[Allegato I]{nis2-directive}. 
        Per la loro rilevanza strategica, questi soggetti sono sottoposti a controlli regolari \textit{ex ante} da parte 
        delle autorità competenti, che verificano l'adempimento in materia NIS2\footcite[art. 32]{nis2-directive}.
    \item \textbf{Soggetti importanti}, ossia gli operatori che, pur non raggiungendo lo stesso livello di 
        criticità sistemica, svolgono comunque un ruolo rilevante per la continuità dei servizi e della catena 
        del valore (tra cui servizi postali, produzione alimentare, gestione dei rifiuti)\footcite[Allegato II]{nis2-directive}. 
        Per essi, la vigilanza avviene prevalentemente \textit{ex post}, cioè a seguito di incidenti o segnalazioni di non conformità\footcite[art. 33]{nis2-directive}.
\end{itemize}
Questa distinzione consente di graduare gli obblighi, i meccanismi di vigilanza e le eventuali sanzioni in 
base alla posizione ricoperta dalle organizzazioni e al loro impatto potenziale sulla resilienza complessiva dell’Unione.

\begin{table}[!ht]
    \centering
    \caption{Settori e soggetti coinvolti da NIS2: dettaglio e classificazione\footcite{tabella-acn}}
    \label{tab:nis2-settori}
    \definecolor{myRed}{HTML}{B71C1C} % dark red
    \definecolor{myOrange}{HTML}{F9A825} % dark yellow
    \definecolor{myGreen}{HTML}{1B5E20} % dark green
    \definecolor{myBlue}{HTML}{0D47A1} % dark blue
    \setlength{\tabcolsep}{6pt}
    \renewcommand{\arraystretch}{1.5}
    \resizebox{\textwidth}{!}{{\large
    \begin{tabular}{|l|l|l|l|l|}
    \hline
        \rowcolor{black}
        \color{white}\textbf{Settore} & \color{white}\textbf{Dettaglio} & \color{white}\textbf{Grandi imprese} & \color{white}\textbf{Medie imprese} & \color{white}\textbf{Piccole e micro imprese} \\ \hline
        \rowcolor{gray!20}\multicolumn{5}{|c|}{\textbf{SETTORI ALTAMENTE CRITICI}} \\ \hline
        Energia & 19 tipologie di soggetto & \cellcolor{myRed}\color{white}Essenziali & \cellcolor{myOrange}\color{white}Importanti & \cellcolor{myGreen}\color{white}Fuori ambito \\ \hline
        Trasporti & 10 tipologie di soggetto & \cellcolor{myRed}\color{white}Essenziali & \cellcolor{myOrange}\color{white}Importanti & \cellcolor{myGreen}\color{white}Fuori ambito \\ \hline
        Settore bancario & DORA \textit{lex specialis} & \cellcolor{myRed}\color{white}Essenziali & \cellcolor{myOrange}\color{white}Importanti & \cellcolor{myGreen}\color{white}Fuori ambito \\ \hline
        Infrastrutture dei mercati finanziari & DORA \textit{lex specialis} & \cellcolor{myRed}\color{white}Essenziali & \cellcolor{myOrange}\color{white}Importanti & \cellcolor{myGreen}\color{white}Fuori ambito  \\ \hline
        Settore sanitario & 5 tipologie di soggetto & \cellcolor{myRed}\color{white}Essenziali & \cellcolor{myOrange}\color{white}Importanti & \cellcolor{myGreen}\color{white}Fuori ambito \\ \hline
        Acqua potabile & 1 tipologia di soggetto & \cellcolor{myRed}\color{white}Essenziali & \cellcolor{myOrange}\color{white}Importanti & \cellcolor{myGreen}\color{white}Fuori ambito \\ \hline
        Acque reflue & 1 tipologia di soggetto & \cellcolor{myRed}\color{white}Essenziali & \cellcolor{myOrange}\color{white}Importanti & \cellcolor{myGreen}\color{white}Fuori ambito \\ \hline
        Infrastrutture digitali & 9 tipologie di soggetto & \cellcolor{myRed}\color{white}Essenziali & Essenziali/Importanti/Fuori ambito & Essenziali/Importanti/Fuori ambito \\ \hline
        Gestione dei servizi TIC (B2B) & 2 tipologie di soggetto & \cellcolor{myRed}\color{white}Essenziali & \cellcolor{myOrange}\color{white}Importanti & \cellcolor{myGreen}\color{white}Fuori ambito \\ \hline
        Spazio & 1 tipologia di soggetto & \cellcolor{myRed}\color{white}Essenziali & \cellcolor{myOrange}\color{white}Importanti & \cellcolor{myGreen}\color{white}Fuori ambito \\ \hline
        \rowcolor{gray!20}\multicolumn{5}{|c|}{\textbf{SETTORI CRITICI}} \\ \hline
        Servizi postali e di corriere & 1 tipologia di soggetto & \cellcolor{myOrange}\color{white}Importanti & \cellcolor{myOrange}\color{white}Importanti & \cellcolor{myGreen}\color{white}Fuori ambito \\ \hline
        Gestione dei rifiuti & 1 tipologia di soggetto & \cellcolor{myOrange}\color{white}Importanti & \cellcolor{myOrange}\color{white}Importanti & \cellcolor{myGreen}\color{white}Fuori ambito \\ \hline
        Fabbricazione, produzione e distribuzione di sostanze chimiche & 1 tipologia di soggetto & \cellcolor{myOrange}\color{white}Importanti & \cellcolor{myOrange}\color{white}Importanti & \cellcolor{myGreen}\color{white}Fuori ambito \\ \hline
        Produzione, trasformazione e distribuzione di alimenti & 1 tipologia di soggetto & \cellcolor{myOrange}\color{white}Importanti & \cellcolor{myOrange}\color{white}Importanti & \cellcolor{myGreen}\color{white}Fuori ambito \\ \hline
        Fabbricazione & 6 tipologie di soggetto & \cellcolor{myOrange}\color{white}Importanti & \cellcolor{myOrange}\color{white}Importanti & \cellcolor{myGreen}\color{white}Fuori ambito \\ \hline
        Fornitori di servizi digitali & 4 tipologie di soggetto & Importanti/Fuori ambito & Importanti/Fuori ambito & Importanti/Fuori ambito \\ \hline
        Ricerca & 2 tipologie di soggetto & \cellcolor{myOrange}\color{white}Importanti & \cellcolor{myOrange}\color{white}Importanti & \cellcolor{myGreen}\color{white}Fuori ambito \\ \hline
        \rowcolor{gray!20}\multicolumn{5}{|c|}{\textbf{ULTERIORI TIPOLOGIE DI SOGGETTI}} \\ \hline
        Pubblica Amministrazione centrale & 4 categorie di PA & \cellcolor{myRed}\color{white}Essenziali & \cellcolor{myRed}\color{white}Essenziali & \cellcolor{myRed}\color{white}Essenziali \\ \hline
        Pubblica Amministrazione regionale e locale & 11 categorie di PA & \cellcolor{myOrange}\color{white}Importanti & \cellcolor{myOrange}\color{white}Importanti & \cellcolor{myOrange}\color{white}Importanti \\ \hline
        Ulteriori tipologie di soggetti & 4 tipologie di soggetti & \cellcolor{myBlue}\color{white}Identificazione dell’Autorità & \cellcolor{myBlue}\color{white}Identificazione dell’Autorità & \cellcolor{myBlue}\color{white}Identificazione dell’Autorità \\ \hline
    \end{tabular}
    }}
\end{table}

L’estensione dell’ambito soggettivo comporta conseguenze rilevanti: molte imprese che prima non erano soggette 
a obblighi specifici di cybersicurezza devono ora implementare sistemi strutturati di gestione del rischio, 
rafforzare le proprie capacità di rilevazione e risposta agli incidenti e predisporre procedure di comunicazione 
rapide ed efficaci con le autorità di vigilanza. 


\subsection{Gestione del rischio}
L'elemento cardine su cui si fonda la Direttiva NIS2 è la gestione del rischio
quale base delle pratiche e delle politiche di sicurezza. L'articolo 21 sancisce esplicitamente l'obbligo,
per i soggetti interessati, di adottare misure tecniche, operative e organizzative adeguate e proporzionate al rischio
cui è esposta l'azienda o l'area dell'azienda in questione\footcite[art. 21]{nis2-directive}. Ciò 
significa che alcuni controlli possono applicarsi sempre o solamente in corrispondenza di livelli di rischio
superiori alle soglie fissate internamente dall'impresa, che in ogni caso ha l'onere di documentare e
giustificare le proprie scelte\footcite[art. 2]{regolamento-2690}.

Il Regolamento di esecuzione (UE) 2024/2690, che stabilisce le norme per l'applicazione della direttiva,
specifica ulteriormente i requisiti tecnici e metodologici connessi all’articolo 21, delineando un insieme 
di pratiche di governance del rischio che i soggetti interessati sono tenuti ad applicare. In particolare, 
il regolamento stabilisce che le politiche di sicurezza e le procedure di gestione del rischio siano formalizzate, 
approvate dagli organi direttivi e soggette a revisione periodica, rafforzando l’idea che la gestione del rischio 
costituisca un processo sistematico e continuativo\footcite[Allegato]{regolamento-2690}. Inoltre, 
esse devono essere soggette a revisione almeno ogni due anni o comunque a seguito di incidenti 
significativi, variazioni organizzative o mutamenti del contesto delle minacce in cui si opera.

Infine, ENISA dettaglia una \textit{Risk Management Policy} che preveda indicatori e misure per monitorare maturità e performance, 
liste di policy e una chiara attribuzione di ruoli/responsabilità, in modo che gli esiti di compliance e review alimentino 
coerentemente il ciclo decisionale del management (risk acceptance documentata, piani correttivi, riesami almeno annuali o a 
seguito di incidenti significativi).\footcite[pp. 9--25]{enisa-implementing-guidance}.


\subsection{Misure di cybersicurezza}
NIS2 introduce un quadro prescrittivo di misure minime che tutte le entità rientranti nel campo di 
applicazione devono adottare. L’articolo 21, paragrafo 2, elenca una serie di ambiti specifici 
che costituiscono il nucleo essenziale della sicurezza informatica richiesta: politiche di analisi e gestione 
del rischio, capacità di gestione degli incidenti, misure di continuità operativa e piani di disaster recovery, 
sicurezza della supply chain e delle relazioni con i fornitori (\textit{Third Party Risk Management}), 
requisiti di sicurezza per lo sviluppo e la manutenzione dei sistemi, testing e auditing periodici, 
impiego della crittografia e di adeguati sistemi di controllo degli accessi, formazione del personale 
e pratiche di \textit{cyber hygiene}\footcite[art. 21]{nis2-directive}. 

Il regolamento di esecuzione, come nel caso precedente, precisa i requisiti tecnici e metodologici 
delle misure di cui all’articolo 21. Esso stabilisce che le politiche di sicurezza devono essere formalizzate 
in documenti approvati dagli organi direttivi, rese disponibili all’interno dell’organizzazione e sottoposte a revisione periodica o a riesame 
in occasione di incidenti significativi e di mutamenti rilevanti del contesto di rischio\footcite[art. 3]{regolamento-2690}.

Le misure minime di sicurezza devono comprendere
la definizione di ruoli e responsabilità interne, la gestione della supply chain e l’inclusione di clausole 
contrattuali con i fornitori, la gestione delle vulnerabilità, la protezione degli asset lungo il loro ciclo 
di vita, procedure di autenticazione forte e controllo degli accessi, misure di sicurezza fisica, protezione 
dei dati e delle comunicazioni, monitoraggio degli eventi di sicurezza e procedure di risposta agli incidenti 
e di ripristino\footcite[Allegato]{regolamento-2690}.

La specificazione puntuale dei controlli che i soggetti devono implementare è lasciata ai singoli
Stati membri, secondo quanto previsto dagli ordinamenti nazionali.


\subsection{Obblighi di notifica}
Un'ulteriore misura introdotta da NIS2 riguarda il regime dettagliato in materia di 
notifica degli incidenti di sicurezza. L'articolo 23 disciplina le modalità e le tempistiche di 
segnalazione alle autorità competenti e ai CSIRT nazionali. La direttiva prevede che entro 
24 ore dall’avvenuta rilevazione di un incidente debba essere trasmessa una comunicazione preliminare, 
seguita entro 72 ore da una notifica completa contenente una prima valutazione 
dell’impatto e delle misure adottate\footcite[art. 23]{nis2-directive}. Entro un mese dall’incidente, l’entità 
è tenuta a fornire una relazione finale che includa analisi approfondita delle cause, impatto sui servizi 
e misure correttive adottate. La normativa contempla inoltre la possibilità di relazioni intermedie, 
quando la gestione dell’incidente lo renda necessario.


\subsection{Governance aziendale e responsabilità}
Un aspetto innovativo e centrale della direttiva è l’introduzione di responsabilità 
specifiche per gli organi di gestione delle entità incluse nel campo di applicazione. 
L’articolo 20 sancisce infatti che i membri degli organi di amministrazione, di direzione o di 
supervisione siano direttamente responsabili dell’approvazione delle misure di gestione del rischio 
e della supervisione della loro attuazione\footcite[art. 20]{nis2-directive}. 
La direttiva richiede inoltre che gli organi di gestione garantiscano che i dipendenti ricevano 
formazione adeguata e regolare in materia di cybersicurezza, al fine di conoscere ed applicare 
in maniera consapevole le misure predisposte. 

Il regolamento di esecuzione rafforza ulteriormente tale approccio, richiedendo
che le responsabilità siano chiaramente attribuite, mediante la creazione di un framework aziendale,
così da garantire una gerarchia aziendale ben definita e favorire la supervisione
attiva da parte dei livelli superiori\footcite[Allegato]{regolamento-2690}.

La responsabilità dei vertici aziendali si estende anche alle conseguenze in caso di inadempienza. 
Sono infatti previste sanzioni amministrative pecuniarie e, nei casi più gravi, 
misure che possono includere la sospensione temporanea dei dirigenti responsabili
dai rispettivi incarichi. 



% 1.3 %%%%%%%%%%%%%%%%%%%%%%%%%%%%%%%%%%%%%%%%%%%%%%%%%%%%%%%%%%%%%%%%%%%%%%%%
\section{Direttiva NIS2 in Italia: ACN e CSIRT Italia}
Il recepimento della Direttiva NIS2 in Italia è avvenuto con il 
\textit{Decreto Legislativo del 4 settembre 2024, n. 138}, entrato in vigore nell'ottobre 2024\footcite{dlgs138}.
Il decreto riprende le principali disposizioni della direttiva, precisa le modalità di applicazione a livello nazionale
e specifica l'attribuzione dei compiti di vigilanza e ricezione delle notifiche all'ACN e al CSIRT Italia.

L’ACN (Agenzia per la Cybersicurezza Nazionale), istituita nel 2021 come agenzia governativa con il compito di proteggere gli interessi nazionali 
in materia di cybersicurezza, ricopre un ruolo centrale nella supervisione delle entità NIS2 essenziali e importanti. 
I suoi compiti principali riguardano la vigilanza sul rispetto degli obblighi di gestione del rischio, 
la ricezione delle notifiche di incidenti significativi, l’adozione di provvedimenti correttivi e 
la comminazione delle sanzioni\footcite{dlgs138}. 

Inoltre, all’ACN è attribuita la facoltà di emanare linee guida vincolanti e istruzioni operative. Tra i documenti
più importanti prodotti dall'agenzia, vi sono la Determinazione ACN 164179\footcite{det-acn} e i suoi allegati 1 e 2
riportanti rispettivamente l'elenco dettagliato delle misure di sicurezza obbligatorie rispettivamente
per i soggetti importanti\footcite{det-acn-allegato1} ed essenziali\footcite{det-acn-allegato2}. Gli allegati sono resi disponibili anche
in formato xls, per facilitarne l'adozione per le aziende.

Accanto all’ACN opera il CSIRT (Computer Security Incident Response Team) Italia, 
responsabile della gestione tecnica delle segnalazioni e del supporto operativo alle imprese. 
Il CSIRT rappresenta il punto di raccolta e analisi delle notifiche 
di incidente trasmesse dai soggetti NIS2 e si occupa di fornire riscontri, raccomandazioni e 
supporto tecnico in risposta alle segnalazioni ricevute\footcite[art. 24]{dlgs138}.
Il suo compito è quindi cruciale per quanto riguarda la condivisione di informazioni su minacce emergenti, 
oltre che per la cooperazione con i CSIRT degli altri Stati membri e con ENISA.

Dal punto di vista delle imprese, l’ACN e il CSIRT costituiscono gli interlocutori designati 
per le questioni legate alla conformità NIS2 sul territorio italiano.
La normativa nazionale prevede, a tal fine, l’obbligo per ciascun soggetto 
di designare una persona fisica avente il ruolo di punto di contatto interno per la cybersicurezza,
con la responsabilità di mantenere un canale di comunicazione con le autorità\footcite[art. 21]{dlgs138}. 
Tale figura svolge una funzione anche strategica, assicurando che l’azienda sia pronta a 
segnalare incidenti e a mettere in atto le disposizioni ricevute, gestendo il flusso 
di informazioni provenienti dall’ACN e dal CSIRT.



% 1.4 %%%%%%%%%%%%%%%%%%%%%%%%%%%%%%%%%%%%%%%%%%%%%%%%%%%%%%%%%%%%%%%%%%%%%%%%
\section{Implicazioni organizzative della Direttiva NIS2}
In questa sezione si analizzano le implicazioni organizzative dell’adozione della Direttiva NIS2, le sfide dal punto di 
vista della governance aziendale, della gestione del rischio e dell’allocazione delle risorse, nonché gli impatti 
sui processi interni e sulla cultura della sicurezza. Per l'analisi si segue la struttura individuata
da un report condotto da Deloitte nel 2023 sulla normativa\footcite{deloitte-nis2}, organizzato in quattro sezioni:
\begin{enumerate}
    \item responsabilità del rischio;
    \item misure di sicurezza informatica;
    \item sicurezza della \textit{supply chain};
    \item notifica degli incidenti.
\end{enumerate}
Inoltre, si aggiungono considerazioni sulle pratiche di formazione e igiene informatica e 
una riflessione sulle principali sfide che le organizzazioni si trovano ad affrontare
per garantire la propria conformità alla direttiva. 


\subsection{Governance e responsabilità nella gestione del rischio}
Il grande impatto di NIS2 sulla realtà aziendale è dovuto in primo luogo alla formale attribuzione della responsabilità
della sicurezza agli organi dirigenziali dell'azienda. Gli esperti sottolineano come, in realtà, il CdA (Consiglio di Amministrazione)
possedesse già questa responsabilità, dovendo rispondere direttamente di ciò che accade all'interno dell'organizzazione, ma 
la formalizzazione di ciò ha catturato l'attenzione dei manager per due motivi principali, secondo Rivieri:
\begin{itemize}
    \item L'attribuzione diretta della responsabilità al board implica la non delegabilità della stessa (non specificata dalla direttiva);
    dunque, il CISO, figura di carattere tecnico, rimane responsabile solamente dell'implementazione delle misure necessarie, mentre
    gli amministratori rispondono anche legalmente degli eventi di sicurezza informatica e perciò sono responsabili dell'accettazione dei rischi correlati.
    \item La formalizzazione pubblica di tale responsabilità costituisce una leva significativa per il CISO nei confronti del top management; in assenza di 
    adeguate competenze per valutare il rischio informatico, i vertici tendono ad affidarsi con maggiore sistematicità alle figure tecniche, anche per 
    ridurre l’esposizione a sanzioni, portando a una maggiore rilevanza del CISO nel processo decisionale.
\end{itemize}

Il CISO, in questo contesto, rappresenta quindi il punto di incontro tra compliance normativa e attuazione tecnica. Secondo ENISA, per 
garantire efficacia tale figura deve disporre di risorse dedicate e autonomia decisionale, 
nonché della possibilità di riferire direttamente al board\footcite[p. 23]{enisa-cybersecurity-roles}. 
Il CISO deve inoltre coordinare un team multidisciplinare, cooperando con figure come \textit{Risk Manager}, \textit{Incident and Threat Manager}, 
\textit{Compliance Officer} e \textit{Supply Chain Security Specialist}, come suggerito da ENISA in conformità con l'ECSF 
(European Cybersecurity Skills Framework). Dal punto di vista delle competenze richieste, 
quindi, deve combinare abilità tecniche (gestione delle vulnerabilità, incident response) con competenze manageriali e comunicative, 
indispensabili per interagire con il board e guidare la cultura organizzativa della sicurezza. 
Scompare la figura del manager IT come tecnico verticale: oggi le aziende necessitano di leader con competenze anche strategiche, 
chiamati ad identificare i rischi aziendali e a supportare decisioni di investimento
e governance nell'area della cybersecurity\footcite[pp. 32--34]{kouns-ciso}.

Infine, oltre all’approvazione e alla supervisione delle misure, le linee guida ENISA precisano meccanismi di compliance monitoring e 
di independent review a supporto del board; si ritengono necessari report indipendenti e standardizzati con metriche chiave per l'identificazione dei rischi,
da presentare almeno annualmente agli organi di gestione per garantire un processo decisionale adeguatamente informato e concretizzare l'accountability
richiesta da NIS2\footcite[pp. 21--25]{enisa-implementing-guidance}.


\subsection{Misure di sicurezza}
La conformità alla Direttiva NIS2 richiede alle imprese l’adozione di un insieme di misure 
organizzative e procedurali che integrino la cybersicurezza nei processi aziendali.

In primo luogo, la realizzazione di politiche di cybersicurezza, formalmente approvate dal board 
e inserite nel sistema di gestione del rischio aziendale, come parte integrante 
della strategia organizzativa\footcite[Allegato]{regolamento-2690}.
In quest'ottica, sono essenziali il mantenimento di registri di catalogo dei rischi, la definizione di indicatori 
di performance relativi alla sicurezza, la presentazione di report periodici agli organi direttivi 
e la revisione periodica delle policy, anche in seguito a incidenti significativi.

Per raggiungere efficacemente i suoi obiettivi, NIS2 richiede un approccio trasversale, 
che coinvolga non solo l’area IT, ma anche i dipartimenti legale, compliance, risorse umane 
e approvvigionamento, oltre alle funzioni operative e al top management.
Le linee guida all'implementazione della direttiva sottolineano come tale coinvolgimento si traduca in attività concrete:
l’ufficio legale e la funzione compliance supportano la produzione della documentazione e 
delle notifiche alle autorità; le risorse umane integrano la 
formazione in materia di cybersicurezza nei programmi di formazione del personale e del board; 
la funzione di approvvigionamento deve inserire requisiti di sicurezza e clausole vincolanti 
nei contratti con i fornitori; il top management, infine, è chiamato a supervisionare e approvare 
tali misure, assumendo responsabilità dirette\footcite{ecso-whitepaper-nis2}\footcite{enisa-implementing-guidance}.

Sul piano operativo, NIS2 e le linee guida ENISA richiedono la predisposizione di procedure strutturate:  
\begin{itemize}
    \item simulazioni periodiche di risposta agli incidenti, documentate e analizzate in termini di \textit{lessons learned};  
    \item audit interni regolari, finalizzati a verificare l’efficacia delle misure di sicurezza e a pianificare il miglioramento; 
    \item testing sistematico delle difese aziendali (penetration test e vulnerability assessment, attività di competenza del CISO e del suo team di cybersicurezza), 
    con l'obbligo di tracciarne i risultati, le azioni correttive intraprese e gli indicatori di performance associati; 
    \item revisione delle misure e aggiornamento delle procedure a seguito di incidenti rilevanti o 
    evoluzioni del contesto operativo.
\end{itemize}

Tutte queste attività devono essere integrate in un ciclo continuo di revisione e miglioramento, 
con il supporto di metriche e indicatori delle performance di sicurezza, e sostenute da un’adeguata documentazione, 
necessaria per dimostrare la conformità alle autorità competenti, in linea con il principio di accountability.


\subsection{Sicurezza della supply chain}
La \textit{supply chain} rappresenta uno dei punti più vulnerabili per la resilienza complessiva 
delle organizzazioni\footcite[pp. 22--25]{ecso-whitepaper-nis2}. 
Perciò, la gestione dei rischi della catena di fornitura richiede un'attenta
mappatura dei fornitori critici, analizzando non solo gli interlocutori diretti ma anche quelli 
di secondo e terzo livello, che possono incidere sulla continuità operativa dei servizi essenziali. 

Le raccomandazioni di ECSO ed ENISA sottolineano l’importanza 
di integrare clausole contrattuali vincolanti nei rapporti con i fornitori, così da garantire 
che anche soggetti esterni adottino misure equivalenti, come richiesto dalla direttiva\footcite[p. 2]{ecso-whitepaper-nis2}\footcite[p. 58]{enisa-implementing-guidance}. 
Ciò implica una stretta collaborazione tra la funzione di approvvigionamento, il CISO e
l’area legale per inserire nei contratti obblighi di notifica degli incidenti, 
standard di sicurezza minimi e diritti di audit da parte dell’azienda cliente.  

Alle clausole si devono affiancare meccanismi di monitoraggio continuo, che comprendano la valutazione periodica 
dei fornitori, la richiesta di attestazioni o certificazioni e, ove necessario, la conduzione di audit.
NIS2 infatti non consente più di limitarsi a un approccio passivo, ma richiede di strutturare processi di
supervisione attivi, sistematici e documentati\footcite[pp. 57--63]{enisa-implementing-guidance}.

Tali misure di sicurezza della supply chain introdotte da NIS2 possono rappresentare anche strumenti per ridurre i costi 
di coordinamento e di transazione tra imprese\footcite{testo-som}. L’approccio regolatorio tende infatti a internalizzare i rischi, in questo caso,
informatici della filiera, promuovendo clausole contrattuali, audit congiunti e standard comuni di sicurezza. 
Ciò rafforza la coerenza complessiva della catena, riducendo le asimmetrie che ostacolano il coordinamento e fornendo
un approccio operativo comune, in linea con l'obiettivo di uniformità perseguito dalla direttiva.

Infine, un'interessante riflessione sulla \textit{supply chain security} riguarda la potenziale nascita di un nuovo scenario competitivo: 
la sicurezza della filiera non è infatti solamente una questione di compliance, 
ma le imprese che dimostrano di aver implementato misure di sicurezza valide possono rafforzare la fiducia degli stakeholder e 
proporsi come partner affidabili sul mercato, guadagnando un vantaggio competitivo non indifferente rispetto ai rivali.


\subsection{Notifica degli incidenti}
Un importante elemento che, come verrà approfondito nel Capitolo \ref{cap:confronto}, differenzia NIS2 dagli standard internazionali
è la volontà del legislatore europeo di promuovere il processo di condivisione delle informazioni relative agli
eventi di cybersicurezza, in modo da diffondere la conoscenza in materia e contribuire all'innovazione.
Perciò, NIS2 pone grande attenzione in merito agli obblighi di notifica degli incidenti, richiedendo alle imprese
l’adozione di una timeline operativa precisa per l'\textit{incident response}, oltre al conseguente 
obbligo di testare periodicamente tali misure\footcite[pp.~18--20]{enisa-cybersecurity-roles}.

La normativa fissa precisi KPI in caso di incidenti significativi per gli obblighi di notifica
e la comunicazione con le autorità di settore. 
Le organizzazioni hanno ventiquattro ore di tempo per fornire una prima segnalazione di allerta ed entro
settantadue ore deve seguire una comunicazione ufficiale di risposta all’incidente.
Dopo un mese, deve essere emesso un rapporto finale sull’incidente oppure, se l’indagine è ancora in corso, 
il CSIRT può rilasciare un rapporto intermedio. In questo modo si favorisce la collaborazione tra le aziende e 
le agenzie governative deputate al controllo della cybersicurezza, nonché l'assistenza in caso di emergenze e la gestione
coordinata delle stesse qualora interessino al tempo stesso più soggetti.

Dal punto di vista organizzativo, a tal fine ogni impresa deve designare un punto di contatto interno per la cybersicurezza, 
responsabile dei rapporti con l’ACN e con il CSIRT nazionale\footcite[art. 21]{dlgs138}. 
Questa figura, che deve essere una persona fisica (nel caso italiano, registrata presso l'ACN), ha il compito di garantire la trasmissione 
delle notifiche di incidente, i rapporti con le autorità e la diffusione interna delle 
istruzioni ricevute. Ciò comporta un rafforzamento del ruolo di Chief Information Security Officer (CISO), come spiegato da Matteo Herin
nell'intervista realizzata per questo elaborato, il quale assorbe tali funzioni tra i suoi compiti.


\subsection{Formazione e cultura della sicurezza informatica}
Un altro pilastro dell’adempimento NIS2 è rappresentato dalla formazione e dalla diffusione 
di una cultura della cybersicurezza a tutti i livelli aziendali. 
La direttiva stabilisce che i membri degli organi di gestione siano direttamente responsabili 
dell’approvazione delle misure di sicurezza e ne supervisionino l’attuazione\footcite[art. 20]{nis2-directive}, 
ma per svolgere efficacemente tali compiti devono possedere una preparazione adeguata. 
La formazione del board è dunque un requisito essenziale per poter garantire 
una gestione consapevole dei rischi e per assumersi l’accountability richiesta dal quadro normativo.

Il tema non riguarda tuttavia soltanto i vertici. 
La conformità alla direttiva passa anche dalla capacità di diffondere tra tutti i dipendenti pratiche 
di \textit{cyber hygiene}, come l’uso sicuro delle credenziali, la gestione corretta delle 
informazioni aziendali e la segnalazione tempestiva di anomalie\footcite[pp. 99--103]{enisa-implementing-guidance}. 
In questo senso, i programmi di formazione periodica assumono la funzione di ridurre il fattore 
di rischio umano, spesso l’anello più debole della catena di sicurezza. 



\subsection{Sfide per le organizzazioni}
Come visto in questa sezione, l'introduzione di NIS2 nel panorama normativo richiede ingenti misure alle aziende
per adattare i propri processi e la propria organizzazione alle disposizioni legislative. Ciò comporta naturalmente
un notevole impiego di risorse, tema centrale data la larghissima portata della direttiva, che coinvolge soggetti
profondamente diversi nella struttura, nelle dimensioni e quindi anche nelle capacità di investimento.

Appare chiaro che le entità più grandi e strutturate abbiano meno difficoltà nell'adozione di NIS2 rispetto alle 
aziende che per la prima volta si trovano a dover dimostrare accountability in materia di cybersicurezza.
Le aziende di dimensioni maggiori e quelle che rientravano nella normativa NIS precedente o già soggette ad altre
misure, come DORA per l'area economico-finanziaria, dispongono già di una struttura avanzata per la gestione
della sicurezza informatica. Per questi soggetti, NIS2 rappresenta solamente un'integrazione necessaria rispetto
ai sistemi esistenti, con novità da introdurre soprattutto per la gestione delle terze parti, degli obblighi di
notifica e delle responsabilità formali del CdA, ma senza alcun cambiamento sostanziale e profondamente impattante\footcite[pp. 2--5]{ecso-whitepaper-nis2}.

In questo contesto è importante notare come la capacità di integrare la sicurezza informatica nei processi decisionali 
e nelle strategie di business rappresenti un fattore di vantaggio competitivo, aumentando la fiducia del mercato 
e la continuità operativa e generando barriere all’entrata per i concorrenti meno maturi dal punto di vista organizzativo\footcite{testo-som}.
Tra essi vi sono i soggetti di dimensioni inferiori o coloro i quali per la prima volta si trovano a dover
implementare misure di cybersicurezza strutturate. In questa categoria rientrano anche le aziende non classificate NIS2, 
ma comunque facenti parte della catena produttiva di soggetti NIS2, elemento che costringe le prime a dotarsi dei 
sistemi richiesti dalla direttiva.
Tali organizzazioni si trovano a dover sostenere ingenti investimenti nell'ambito cyber, necessitando di sistemi
nuovi da progettare ad hoc, di una revisione profonda dei processi operativi e di figure esperte che 
abbiano le competenze per implementare tutto il necessario, con un conseguente impatto elevato sulle proprie risorse\footcite[pp. 2--5]{ecso-whitepaper-nis2}.

Le aziende che si trovano in questa situazione, come si vedrà nei capitoli successivi, possono decidere di realizzare
i controlli richiesti strutturando direttamente la sicurezza informatica sulle disposizioni NIS2 oppure adottando un 
framework che consenta loro di progettare una base solida per la cybersicurezza, su cui integrare gli adempimenti mancanti. 
In entrambe le casistiche, tuttavia, la complessità organizzativa
rimane elevata e si richiede un'accurata pianificazione delle strategie di business e del budget per includere l'integrazione
di NIS2 nei processi.
